\section{Financial Education Pilot in Elementary and Middle Schools}\label{sec:pilot}

The National Strategy for Financial Education (\textit{Estratégia Nacional de Educação Financeira} -- ENEF) was established by the Brazilian Federal Government in 2010 to promote financial education and empower individuals to make autonomous financial decisions.\footnote{ENEF was instituted by Federal Decree 7.397 of December 22, 2010. It brings together the Ministry of Education, the Central Bank of Brazil, the Securities and Exchange Commission of Brazil (CVM), the Ministry of Finance, the Superintendence of Private Insurance (SUSEP), and the Superintendence of Complementary Pension (PREVIC).} Between August 2010 and December 2011, ENEF implemented a pilot financial education program in public high schools across six states -- São Paulo, Rio de Janeiro, Ceará, Tocantins, Minas Gerais, and the Federal District. The program's impact was assessed through a randomized controlled trial involving 891 schools and approximately 20,000 students, and generated significant improvements in students' financial knowledge, saving intentions, financial autonomy, and participation in household financial decisions \citep{Bruhn2016}. Building on these results, in 2014 ENEF commissioned the Brazilian National Committee on Financial Education (\textit{Comitê Nacional de Educação Financeira} -- CONEF) to extend the program to elementary and middle school grades. The pilot of this new program is the object of our evaluation.

The pilot was implemented during the 2015 school year in municipal schools of Joinville and Manaus, whose respective Municipal Secretariats of Education voluntarily joined the initiative. The two cities represent distinct socioeconomic contexts: Manaus, the capital of Amazonas, lies in the northern part of the country, while Joinville, a municipality of Santa Catarina, is located in the more affluent South. By 2016, Joinville exhibited higher per capita income (USD 14,000 vs.~USD 11,000), a higher Municipal Human Development Index (0.809 vs.~0.737), higher school attendance rates among 6- to 14-year-olds (97\% vs.~94\%), and stronger performance on the Brazilian Education Development Index (IDEB) at both elementary and middle school levels. While these differences introduce regional variation in the environment in which the program was delivered, they also extend the external validity of our findings to two economic and educational contexts that are representative of much of Brazil's public-school network.

\subsection{Teaching Materials}\label{subsec:materials}

The program rests on a pair of teaching collections published in 2014 by CONEF -- a student book (\textit{Livro do Aluno}) and a teacher book (\textit{Livro do Professor}), with one volume per grade from the 1\textsuperscript{st} to the 9\textsuperscript{th} year, totaling 18 volumes.\footnote{The materials are publicly available through the ENEF portal and were coordinated editorially by the Brazilian stock exchange (BM\&FBOVESPA, now B3). They were validated in 2011 by an inter-institutional pedagogical group comprising representatives of the Ministry of Education, the Central Bank of Brazil, the CVM, the Ministry of Finance, SUSEP, PREVIC, the University of Brasília, the Federal University of Rondônia, the Federal Institute of Education of Ceará, the Federal University of Rio Grande do Sul, and Colégio Pedro II.} The program adopts the OECD's definition of financial education, with emphasis on the ``formation'' vertent -- that is, on cultivating attitudes and behaviors rather than on the provision of product-specific information -- on the grounds that children and adolescents lack direct access to financial products \citep{Lusardi2014}. Its content is organized along two dimensions. The \textit{spatial} dimension moves from the individual to the local, regional, national, and global levels, emphasizing how individual choices reverberate at larger scales. The \textit{temporal} dimension connects past decisions to present circumstances and present decisions to future outcomes, placing the intertemporal nature of financial choice at the center. The curriculum articulates six objectives -- four of them spatial (educating for citizenship; ethical and responsible consumption and saving; autonomy in financial decisions; and training of multipliers) and two temporal (planning at short-, medium-, and long-term horizons; and a culture of prevention) -- which in turn map onto ten learning competencies.

The content is not delivered as a stand-alone subject. Instead, it is designed to be \textit{transversally integrated} into the existing curriculum and is intended to be taught by teachers of any specialty, without requiring prior financial training.\footnote{The program rests explicitly on Morin's notion of \textit{religação dos saberes} (the reconnection of knowledge domains). The teacher books instruct educators to deliver the content ``through their role as citizens rather than through their disciplinary specialty''.} This integration is built into the material itself: the margins of the \textit{Livro do Professor} contain icons that flag, on each page, where the activity connects with Portuguese (reading comprehension, authorship, genre), Mathematics (budgeting, estimation, tabular representation), or Environmental Education, as well as where a formal financial concept or a psychological bias (\textit{armadilha psicológica}) is introduced. We discuss in \autoref{sec:implementation} how this transversality played out in the classroom.

Pedagogically, the program follows two distinct blocks. The \textit{early grades} (1\textsuperscript{st}--4\textsuperscript{th}) follow a project-based approach structured around four thematic axes -- Production and Consumption, Organization, Care, and Planning -- revisited each year with age-appropriate objects. In the 3\textsuperscript{rd} grade, these four axes are organized around a ball, the household's expenses, the school, and Book Day festivities. The \textit{later grades} (5\textsuperscript{th}--9\textsuperscript{th}) follow a narrative-based approach in which students learn concepts, attitudes, and biases by making decisions within a fictional story, whose format itself changes across grades. \autoref{tab:materials} summarizes the format and central content of each of the four volumes delivered to pilot grades. Assessment of learning, in the program's own design, relies on student self-reflection rather than on formal testing; the standardized financial proficiency instrument we describe in \autoref{sec:data_c} is therefore external to the program and is not part of its internal pedagogy.

\begin{table}[h!]
\centering
\fontsize{9}{11}\selectfont
\begin{threeparttable}
\caption{Financial education materials by pilot grade}\label{tab:materials}
\begin{tabular}{p{1.1cm}p{5cm}p{8.2cm}}
\toprule
Grade & Pedagogical format & Central concepts and attitudes \\
\midrule
3\textsuperscript{rd} & Four projects organized around concrete objects of everyday life: a ball (Production and Consumption), the household's expenses (Organization), the school (Care), and Book Day festivities (Planning). & Price composition; environmentally responsible consumption; needs vs.~wants; income and expenses; care for shared public and private spaces; short-term planning and estimation. \\
\addlinespace
5\textsuperscript{th} & Three choose-your-own-adventure stories (\textit{aventuras-solo}) plus a closing group project, centered on environmental sustainability and the ``five R's'' (rethink, refuse, reduce, reuse, recycle). & Waste; cash vs.~credit purchases; value vs.~price; conscious consumption; short-term planning; risk; patrimony; sustainability. Psychological traps introduced: ostentation, immediacy, selective perception, over-confidence. \\
\addlinespace
7\textsuperscript{th} & A single serialized game over six meetings in which groups representing different economic agents negotiate proposals and manage a running budget, with new rules added at each meeting. & Personal budget (income $-$ expenses $=$ balance); investment and the risk-return trade-off; financing and interest; the financial system; patrimony protection; borrowing. Additional traps: framing, sunk costs, loss aversion. \\
\addlinespace
9\textsuperscript{th} & A print publication designed to simulate a news website (\textit{impressite}), with reports, interviews, columns, a three-episode web-series, a forum, and hands-on exercises. & Household budgeting; expense categorization; consumer traps; social entrepreneurship; the Consumer Defense Code; public space, public budget, and taxes. Attitudes emphasized: autonomy, risk perception, discipline, and spending control. \\
\bottomrule
\end{tabular}
\begin{tablenotes}
\fontsize{9}{9}\selectfont
\item Notes: Based on the \textit{Livro do Aluno} and \textit{Livro do Professor}, volumes 3, 5, 7, and 9, CONEF (2014). The four volumes together articulate six common learning objectives (four spatial: education for citizenship, ethical consumption and saving, autonomy in financial decisions, and training of multipliers; two temporal: short-, medium-, and long-term planning and a culture of prevention) and ten corresponding competencies.
\end{tablenotes}
\end{threeparttable}
\end{table}

A direct implication of this design is that the intervention differs across the four pilot grades along three dimensions simultaneously: student age, textbook content, and pedagogical method. We view this as a design feature of the program rather than as something we can isolate empirically: the program is the bundle. Differences in estimated effects across grades therefore cannot be cleanly interpreted as life-cycle or maturation effects, as they also embed differences in the specific content of each volume and in how the content is delivered. We return to this caveat when presenting grade-specific estimates in \autoref{sec:results}.

\subsection{Sample Selection}\label{subsec:sample}

The evaluation was agreed with the implementing partner to include approximately 200 schools, so as to balance statistical power against the operational cost of data collection. We used existing administrative data to define the sampling frame. According to the 2015 Census of Education, Joinville had 72 municipal schools offering at least one of the pilot grades; we included all of them. Manaus had 302 eligible municipal schools; we sampled 129.\footnote{The 129 schools in Manaus are the result of a two-stage procedure. We first excluded 53 schools located in riverside communities, following the guidelines of the Department of Education, due to access constraints. Of the remaining 302 schools offering at least one target grade, we sampled 36 of the 202 schools that only offered elementary grades based on the 2013 IDEB quartiles of their 5\textsuperscript{th}-grade performance; 28 of the 35 schools that only offered middle school grades based on the 2013 IDEB quartiles of their 9\textsuperscript{th}-grade performance; and all 65 schools that offered grades 1 to 9.} In total, 201 schools entered the experiment and were randomly assigned to treatment (101 schools) or control (100 schools).

In agreement with ENEF, we selected grades 3, 5, 7, and 9 to generate evidence on each of the four volumes delivered to pilot grades. The sample comprises 9,204 students in treatment schools and 9,353 students in control schools. To account for differences in infrastructure and socioeconomic context, we stratified the randomization by municipality and by three school types -- schools offering only elementary grades (1\textsuperscript{st}--5\textsuperscript{th}), schools offering only middle school grades (6\textsuperscript{th}--9\textsuperscript{th}), and schools offering all elementary and middle school grades (1\textsuperscript{st}--9\textsuperscript{th}) -- generating six strata in total (\autoref{tab:sample}).

An important feature of the design is that, within each school, the program was not delivered to every class of the pilot grades. At the beginning of the 2015 school year, each school principal -- in both treatment and control schools -- selected a single class at each pilot grade to be included in the survey sample. This strategy was adopted to keep the cost of data collection manageable. It introduces, however, a potential source of endogeneity: in treatment schools, the principal's choice may have been influenced by the willingness of individual teachers to adopt the new material, which could in turn correlate with teacher characteristics or classroom composition. We flag this concern upfront because it bears directly on the internal validity of our estimates. We document the degree of compliance with the treatment at the class level in \autoref{sec:implementation}, and we return to the implications of this within-school selection of classes when interpreting our estimates in \autoref{sec:results}.\footnote{Control schools were not exposed to the financial education material, and principals in control schools had no reason to select classes on characteristics that would correlate with the intervention. As we show in \autoref{sec:data_c}, the observable characteristics of students in the treatment and control groups are well balanced in the overall sample, which suggests that the selection of classes within treated schools, though non-random, did not produce sizable systematic differences along observables in our estimation sample.}
